%% schriftsatz.latex — Anpassbares Pandoc-Template für familienrechtliche Schriftsätze
%%
%% Pandoc-Variablen (via -V oder YAML-Frontmatter):
%%   mainfont      Schriftart             (Standard: Arial)
%%   fontsize      Schriftgröße           (Standard: 11pt)
%%   papersize     Papierformat           (Standard: a4)
%%   geometry      Seitenränder           (Standard: margin=2.5cm)
%%   linestretch   Zeilenabstand          (Standard: 1.2)
%%   lang          Sprache                (Standard: de-DE)
%%   briefkopf     Vorbereiteter LaTeX-Briefkopf (von generate-pdf.py erzeugt)
%%
%% ─────────────────────────────────────────────────────────────────────────────

\documentclass[
  11pt,
  a4paper,
]{article}

%% ── Sprache ──────────────────────────────────────────────────────────────────
\usepackage[ngerman]{babel}

%% ── Fonts (XeLaTeX) ──────────────────────────────────────────────────────────
\usepackage{fontspec}
\setmainfont{Arial}[
  BoldFont       = Arial Bold,
  ItalicFont     = Arial Italic,
  BoldItalicFont = Arial Bold Italic,
  Ligatures      = TeX,
]
\usepackage{microtype}

%% ── Seitenlayout ─────────────────────────────────────────────────────────────
\usepackage[left=3cm,right=2cm,top=2.5cm,bottom=2.5cm]{geometry}

%% ── Zeilenabstand & Absätze ──────────────────────────────────────────────────
\usepackage{setspace}
\setstretch{1.3}
\usepackage{parskip}
\setlength{\parskip}{6pt}
\setlength{\parindent}{0pt}

%% ── Überschriften ────────────────────────────────────────────────────────────
\usepackage{titlesec}
\titleformat{\section}{\normalfont\large\bfseries}{}{0em}{}
\titleformat{\subsection}{\normalfont\normalsize\bfseries}{}{0em}{}
\titleformat{\subsubsection}{\normalfont\normalsize\bfseries}{}{0em}{}
\titlespacing{\section}{0pt}{14pt}{6pt}
\titlespacing{\subsection}{0pt}{10pt}{4pt}
%% Keine automatische Abschnittsnummerierung
\setcounter{secnumdepth}{0}

%% ── Seitenzahlen ─────────────────────────────────────────────────────────────
\usepackage{fancyhdr}
\usepackage{lastpage}
\pagestyle{fancy}
\fancyhf{}
\fancyfoot[C]{\footnotesize Seite~\thepage~von~\pageref{LastPage}}
\renewcommand{\headrulewidth}{0pt}
\renewcommand{\footrulewidth}{0pt}

%% ── Tabellen ─────────────────────────────────────────────────────────────────
\usepackage{booktabs}
\usepackage{longtable}
\usepackage{array}
\usepackage{calc}
\usepackage{caption}
\captionsetup{font=small, labelfont=bf}
%% Pandoc 3.x setzt \LTcaptype auf "none" — LaTeX braucht diesen Counter:
\newcounter{none}

%% ── Listen ───────────────────────────────────────────────────────────────────
\usepackage{enumitem}
\setlist{noitemsep, topsep=4pt}

%% ── Sonstiges ────────────────────────────────────────────────────────────────
\usepackage{xcolor}
\usepackage[normalem]{ulem}
\usepackage[hidelinks]{hyperref}
\usepackage{needspace}

%% ── Pandoc-Kompatibilität ────────────────────────────────────────────────────
\providecommand{\tightlist}{%
  \setlength{\itemsep}{0pt}\setlength{\parskip}{0pt}}

%% ─────────────────────────────────────────────────────────────────────────────
\begin{document}

%% Briefkopf — von generate-pdf.py aus Markdown-Zeilen vor dem ersten --- erzeugt.
%% Zeilen mit "den \d" werden rechtsbündig gesetzt, Betreff fett+unterstrichen,
%% Aktenzeichen fett, alle anderen Zeilen linksbündig mit explizitem Zeilenumbruch.

\section{Vorbereitung Verhandlung --- Az. 7 F
88/25}\label{vorbereitung-verhandlung-az.-7-f-8825}

\subsection{Verfahrensdaten}\label{verfahrensdaten}

\begin{itemize}
\tightlist
\item
  \textbf{Termin:} 22. Mai 2026, 10:00 Uhr
\item
  \textbf{Gericht:} Amtsgericht Neukirchen, Saal 3
\item
  \textbf{Richterin/Richter:} Richterin Dr.~Lieselotte Kern
\item
  \textbf{Verfahrensbeistand:} Dipl.-Soz. Klaus Hofmann
\item
  \textbf{Gegenanwalt/-anwältin:} RA Petra Schmitt
\end{itemize}

\begin{center}\rule{0.5\linewidth}{0.5pt}\end{center}

\subsection{Kernbotschaften (max. 3)}\label{kernbotschaften-max.-3}

\begin{enumerate}
\def\labelenumi{\arabic{enumi}.}
\tightlist
\item
  Finn geht es gut --- das sagt die KiTa, das sagen die Zahlen, das
  zeigt das gelebte Wechselmodell der letzten sechs Monate.
\item
  Der Antragsgegner war von Anfang an aktiv beteiligt und hat seine
  Rolle als Vater nie vernachlässigt.
\item
  Das paritätische Wechselmodell wurde einvernehmlich vereinbart und
  funktioniert --- eine Änderung liegt nicht im Interesse Finns.
\end{enumerate}

\begin{center}\rule{0.5\linewidth}{0.5pt}\end{center}

\subsection{Punkte aus dem Schriftsatz
(Kurzfassung)}\label{punkte-aus-dem-schriftsatz-kurzfassung}

\subsubsection{Eigene Stärken}\label{eigene-stuxe4rken}

\begin{itemize}
\tightlist
\item
  Stabile, kindgerechte Wohnung seit 01.07.2025; gleicher KiTa-Bezirk,
  1,8 km Entfernung zur AS
\item
  Elternvereinbarung vom 01.09.2025 belegt einvernehmliche Einigung (B1)
\item
  AWO-Beratung am 14.01.2026 gemeinsam besucht ---
  Kooperationsbereitschaft dokumentiert
\end{itemize}

\subsubsection{Fachliche
Bestätigungen}\label{fachliche-bestuxe4tigungen}

\begin{itemize}
\tightlist
\item
  KiTa Regenbogen: Halbjahresbericht Herbst 2025 positiv (B2);
  Entwicklungsgespräch 25.02.2026 bestätigt altersgemäße Entwicklung
  (B5)
\end{itemize}

\subsubsection{Widerlegungen}\label{widerlegungen}

\begin{itemize}
\tightlist
\item
  Behauptete Abwesenheit am 15.10.2025 durch Dienstreise-Belege
  widerlegt (B4): AG war in Neukirchen, Dienstreise war am 16.10.
\item
  Druckbehauptung zur Elternvereinbarung: Entstehungsgeschichte
  dokumentiert (B3, AWO)
\end{itemize}

\begin{center}\rule{0.5\linewidth}{0.5pt}\end{center}

\subsection{NUR MÜNDLICH-Punkte}\label{nur-muxfcndlich-punkte}

→ Siehe \texttt{erwiderung/nur-muendlich.md} für Details und
Redevorschläge.

{\def\LTcaptype{none} % do not increment counter
\begin{longtable}[]{@{}
  >{\raggedright\arraybackslash}p{(\linewidth - 4\tabcolsep) * \real{0.1842}}
  >{\raggedright\arraybackslash}p{(\linewidth - 4\tabcolsep) * \real{0.4211}}
  >{\raggedright\arraybackslash}p{(\linewidth - 4\tabcolsep) * \real{0.3947}}@{}}
\toprule\noalign{}
\begin{minipage}[b]{\linewidth}\raggedright
Punkt
\end{minipage} & \begin{minipage}[b]{\linewidth}\raggedright
Wann vorbringen
\end{minipage} & \begin{minipage}[b]{\linewidth}\raggedright
Redevorschlag
\end{minipage} \\
\midrule\noalign{}
\endhead
\bottomrule\noalign{}
\endlastfoot
Spontanes Lob KiTa-Leiterin & Wenn Richterin fragt nach Einschätzung der
KiTa & „Die KiTa-Leiterin hat mich nach dem Gespräch noch kurz zur Seite
genommen\ldots'' \\
Veränderung bei Übergaben seit Antrag & Wenn Gegenseite mangelnde
Kooperation behauptet & „Seit dem Antrag hat sich die Atmosphäre bei den
Übergaben verändert\ldots'' \\
\end{longtable}
}

\begin{center}\rule{0.5\linewidth}{0.5pt}\end{center}

\subsection{Erwartbare Fragen der
Richterin}\label{erwartbare-fragen-der-richterin}

{\def\LTcaptype{none} % do not increment counter
\begin{longtable}[]{@{}
  >{\raggedright\arraybackslash}p{(\linewidth - 2\tabcolsep) * \real{0.2917}}
  >{\raggedright\arraybackslash}p{(\linewidth - 2\tabcolsep) * \real{0.7083}}@{}}
\toprule\noalign{}
\begin{minipage}[b]{\linewidth}\raggedright
Frage
\end{minipage} & \begin{minipage}[b]{\linewidth}\raggedright
Antwortvorschlag
\end{minipage} \\
\midrule\noalign{}
\endhead
\bottomrule\noalign{}
\endlastfoot
„Wie stellen Sie sich die Betreuung konkret vor?'' & Beibehaltung des
aktuellen Wechselmodells: Mo--Mi Mutter, Mi--Fr Vater, Wochenenden
alternierend. Übergabe über KiTa. \\
„Warum sind Sie gegen das Residenzmodell?'' & Finn kennt beide Haushalte
als gleichwertige Zuhause. Ein Residenzmodell würde eine funktionierende
Situation ohne Anlass verschlechtern. \\
„Wie gehen Sie mit den geschilderten Symptomen um?'' & Die KiTa sieht
keine Auffälligkeiten. Symptome treten laut Protokoll der AS
ausschließlich in ihrem Haushalt auf. Das ist medizinisch einzuordnen,
nicht dem Wechselmodell anzulasten. \\
\end{longtable}
}

\begin{center}\rule{0.5\linewidth}{0.5pt}\end{center}

\subsection{Erwartbare Argumente der
Gegenseite}\label{erwartbare-argumente-der-gegenseite}

{\def\LTcaptype{none} % do not increment counter
\begin{longtable}[]{@{}
  >{\raggedright\arraybackslash}p{(\linewidth - 2\tabcolsep) * \real{0.5000}}
  >{\raggedright\arraybackslash}p{(\linewidth - 2\tabcolsep) * \real{0.5000}}@{}}
\toprule\noalign{}
\begin{minipage}[b]{\linewidth}\raggedright
Argument
\end{minipage} & \begin{minipage}[b]{\linewidth}\raggedright
Reaktion
\end{minipage} \\
\midrule\noalign{}
\endhead
\bottomrule\noalign{}
\endlastfoot
„Das Kind leidet unter dem Wechselmodell'' & KiTa-Bericht und
Entwicklungsgespräch widersprechen dem ausdrücklich. Symptome sind auf
AS-Haushalt beschränkt. \\
„Die Vereinbarung wurde unter Druck unterzeichnet'' &
Entstehungsgeschichte belegt das Gegenteil: mehrwöchige Aushandlung, AWO
einbezogen, beide Seiten in CC bei E-Mails. \\
„Der Vater war beruflich oft abwesend'' & Dienstreisen 2025 gesamt: 10
Tage. Belege liegen vor. In Abwesenheitszeiten war Regelbetreuung durch
AS wie im Modell vorgesehen. \\
\end{longtable}
}

\begin{center}\rule{0.5\linewidth}{0.5pt}\end{center}

\subsection{Dos and Don'ts}\label{dos-and-donts}

\subsubsection{Do:}\label{do}

\begin{itemize}
\tightlist
\item
  Ruhig und sachlich bleiben
\item
  Vom Kind sprechen, nicht vom Konflikt
\item
  Kooperationsbereitschaft zeigen
\item
  Konkrete Beispiele nennen
\item
  Auf Fragen direkt antworten
\end{itemize}

\subsubsection{Don't:}\label{dont}

\begin{itemize}
\tightlist
\item
  Die Mutter angreifen
\item
  Emotional werden
\item
  Unterbrechen
\item
  Rechtfertigungen vorbringen, die nicht gefragt wurden
\item
  Details ausbreiten, die im Schriftsatz nicht stehen
  (Überraschungsgefahr)
\end{itemize}

\begin{center}\rule{0.5\linewidth}{0.5pt}\end{center}

\subsection{Anhang: Verhaltensregeln}\label{anhang-verhaltensregeln}

→ Vollständige Regeln für Gericht, Jugendamt und Verfahrensbeistand:
\texttt{references/verhaltensregeln.md}

\subsubsection{Der wichtigste Satz für das gesamte
Verfahren:}\label{der-wichtigste-satz-fuxfcr-das-gesamte-verfahren}

\subsubsection{Wenn du unter Druck gesetzt oder provoziert
wirst:}\label{wenn-du-unter-druck-gesetzt-oder-provoziert-wirst}

\begin{enumerate}
\def\labelenumi{\arabic{enumi}.}
\tightlist
\item
  Pause --- nicht sofort antworten
\item
  „Das entspricht nicht meiner Wahrnehmung'' statt „Das stimmt nicht''
\item
  Auf Fakten verweisen, nicht auf Gefühle
\item
  Die Richterin entscheidet --- nicht die Gegenanwältin
\end{enumerate}

\subsubsection{Wenn du merkst, dass du emotional
wirst:}\label{wenn-du-merkst-dass-du-emotional-wirst}

Sage nichts. Bitte um einen kurzen Moment. Das ist kein Zeichen von
Schwäche --- es ist Professionalität. Richterinnen kennen das.

\end{document}
